\chapter{Testes e Validação}
\label{cap:testes}

\section{Estratégia de testes}

A validação do Cardly combina testes automatizados no backend, testes manuais integrados mobile-API e evidências visuais exigidas pelo manual acadêmico UNILESTE \cite{junit5_spring,jest_testing_library,postman_docs}. A estratégia prioriza regras de negócio críticas (autenticação, agendamento de cartas, autorização admin).

\section{Testes automatizados -- backend}

Suite JUnit 5 + Spring Boot Test executada via Maven (\texttt{mvn test}) \cite{junit5_spring}. Classes principais:

\begin{table}[htbp]
    \caption{Testes unitários e de integração do backend}
    \label{tab:testes-backend}
    \centering
    \small
    \begin{tabular}{ll}
        \toprule
        \textbf{Classe} & \textbf{Escopo} \\
        \midrule
        AuthServiceTest & Registro, login, hash de senha \\
        UserServiceTest & CRUD usuário, soft delete \\
        DeckServiceTest & Assuntos, visibilidade pública \\
        CardServiceTest & Respostas, intervalos, dueAt \\
        EmailNormalizerTest & Normalização de e-mail \\
        CardlyBackendApplicationTests & Context load Spring \\
        \bottomrule
    \end{tabular}
\end{table}

Ambiente de teste utiliza H2 ou PostgreSQL em container via \texttt{application-test.properties} \cite{docker_postgres}.

\subsection{Exemplo: validação de agendamento}

Testes em \texttt{CardServiceTest} verificam:
\begin{itemize}
    \item Carta nova correta $\rightarrow$ EASY, 5 dias.
    \item Carta nova incorreta $\rightarrow$ MEDIUM, 3 dias.
    \item Bloqueio de resposta antes de \texttt{dueAt}.
    \item Skip reagenda conforme dificuldade atual.
\end{itemize}

\section{Testes manuais -- API}

Postman (ou Insomnia) valida contratos REST \cite{postman_docs}:

\begin{enumerate}
    \item Registrar usuário e obter JWT.
    \item Criar deck e cartas.
    \item Responder carta e consultar novo \texttt{dueAt}.
    \item Tentar acessar \texttt{/api/admin} com role USER (esperado 403).
    \item Buscar decks com payload vazio em \texttt{/search} (paginação).
\end{enumerate}

pgAdmin inspeciona persistência no PostgreSQL \cite{pgadmin_docs}.

\section{Testes manuais -- aplicativo mobile}

Roteiro de validação integrada:

\begin{longtable}{p{1cm} p{5cm} p{7.5cm}}
    \caption{Roteiro de testes manuais mobile} \label{tab:testes-mobile} \\
    \toprule
    \textbf{ID} & \textbf{Cenário} & \textbf{Resultado esperado} \\
    \midrule
    \endfirsthead
    T01 & Registro e login & Token salvo; navegação ao dashboard \\
    T02 & Criar assunto privado & Lista atualizada; toast sucesso \\
    T03 & Adicionar cartas & Cartas visíveis no deck \\
    T04 & Sessão de estudo + flip & Animação; resposta registrada \\
    T05 & Carta agendada futura & Indisponível até dueAt \\
    T06 & Dashboard & Contadores coerentes com API \\
    T07 & Comunidade + clone & Deck copiado na coleção \\
    T08 & Admin (superadmin) & CRUD global funcional \\
    T09 & Excluir conta & Modal confirma; logout automático \\
    \bottomrule
\end{longtable}

\section{Evidências visuais (prints)}

O relatório final deve incluir capturas de tela de diversos cenários reais, conforme requisito institucional. Placeholders abaixo indicam prints pendentes de inserção:

\begin{figure}[htbp]
    \centering
    \fbox{\parbox{0.4\textwidth}{\centering\vspace{2.5cm}Login\vspace{2.5cm}}
    \fbox{\parbox{0.4\textwidth}{\centering\vspace{2.5cm}Estudo + Flip\vspace{2.5cm}}
    \caption{Evidências visuais -- login e sessão de estudo (inserir PNGs reais)}
    \label{fig:prints-estudo}
\end{figure}

\begin{figure}[htbp]
    \centering
    \fbox{\parbox{0.4\textwidth}{\centering\vspace{2.5cm}CRUD Decks\vspace{2.5cm}}
    \fbox{\parbox{0.4\textwidth}{\centering\vspace{2.5cm}Admin Panel\vspace{2.5cm}}
    \caption{Evidências visuais -- CRUD e administração}
    \label{fig:prints-crud}
\end{figure}

Recomenda-se pasta \texttt{docs/report/figures/} com arquivos \texttt{login.png}, \texttt{study.png}, etc.

\section{Testes de front-end (planejado)}

React Native Testing Library \cite{jest_testing_library} pode cobrir:
\begin{itemize}
    \item Renderização de \texttt{FlipCard} nos estados frente/verso.
    \item Exibição condicional de \texttt{AdminScreen}.
    \item Helpers de toast.
\end{itemize}

Cobertura automatizada mobile permanece como melhoria futura; validação atual é predominantemente manual integrada.

\section{CI/CD (opcional)}

GitHub Actions pode executar \texttt{mvn test} a cada push \cite{github_actions}. Pipeline não bloqueia entrega acadêmica, mas reforça estabilidade citada nos critérios de avaliação.

\section{Rastreabilidade requisito-teste}

\begin{table}[htbp]
    \caption{Rastreabilidade entre requisitos e testes}
    \label{tab:rastreabilidade}
    \centering
    \footnotesize
    \begin{tabular}{lll}
        \toprule
        \textbf{RF/RT} & \textbf{Teste} & \textbf{Tipo} \\
        \midrule
        RF01--RF03 & AuthServiceTest, T01, T09 & Auto + manual \\
        RF04--RF06 & DeckServiceTest, T02--T03 & Auto + manual \\
        RF07--RF10 & CardServiceTest, T04--T05 & Auto + manual \\
        RF09 & DashboardController, T06 & Auto + manual \\
        RF15--RF16 & Admin endpoints, T08 & Manual \\
        RT01 & AuthServiceTest, T01 & Auto + manual \\
        RT05 & CardServiceTest search, Postman & Auto + manual \\
        \bottomrule
    \end{tabular}
\end{table}

\section{Critérios de aceite para demonstração}

Para a apresentação prática, considera-se aceito:
\begin{itemize}
    \item Login $\rightarrow$ estudo $\rightarrow$ dashboard sem erro fatal.
    \item API respondendo em produção ou emulador apontando para Railway.
    \item Pelo menos um teste automatizado de carta demonstrado na arguição.
    \item Prints inseridos na versão final deste PDF.
\end{itemize}
